\subsection{Complexes of homomorphisms}

Let (C, d) and (C', d') be two A-complexes. Consider the graded k-module Homgr\_A (C, C') (II, pp. 174–175): for \(n \in \mathbf{Z}\) Homgr\_A (C, C')\_n is the k-module of graded A-linear maps of degree \(n\) from C to C'; in other words, Homgr\_A (C, C') is canonically identified with the A-module

\[
\bigoplus_ {n \in \mathbf {Z}} \prod_ {p \in \mathbf {Z}} \operatorname{Hom} _ {\mathrm{A}} \left(\mathrm{C} _ {p}, \mathrm{C} _ {p + n} ^ {\prime}\right) = \bigoplus_ {n \in \mathbf {Z}} \prod_ {p + q = n} \operatorname{Hom} _ {\mathrm{A}} \left(\mathrm{C} _ {p}, \mathrm{C} ^ {\prime q}\right).
\]

Let us define k-linear maps.

\[
\mathrm{D} _ {n}: \operatorname{Homgr} _ {\mathbf {A}} (\mathrm{C}, \mathrm{C} ^ {\prime}) _ {n} \rightarrow \operatorname{Homgr} _ {\mathbf {A}} (\mathrm{C}, \mathrm{C} ^ {\prime}) _ {n - 1}, \quad n \in \mathbf {Z},
\]

by

\[
\mathbf {D} _ {n} (f) = d ^ {\prime} \circ f - (- 1) ^ {n} f \circ d;\tag{1}
\]

we have

\[
\begin{array}{r l} \mathrm{D} _ {n - 1} \circ \mathrm{D} _ {n} (f) = & \mathrm{D} _ {n - 1} \left(d ^ {\prime} \circ f - (- 1) ^ {n} f \circ d\right) = d ^ {\prime} \circ d ^ {\prime} \circ f - (- 1) ^ {n} d ^ {\prime} \circ f \circ d \\ & - (- 1) ^ {n - 1} d ^ {\prime} \circ f \circ d - f \circ d \circ d = 0. \end{array}
\]

Then (Homgr \(_{A}\) (C, C'), D) is a complex of \(k\) -modules called the complex of homomorphisms from C to C'.

For example, \(\operatorname{Homgr}_{\mathbf{A}}(\mathbf{A}, \mathbf{C}')\) is canonically identified with \(\mathbf{C}'\). Let us also note that, for every \(n \in \mathbf{Z}\) , we have \(\operatorname{Homgr}_{\mathbf{A}}(\mathbf{C}, \mathbf{C}')\) ( \(n\) ) = \(\operatorname{Homgr}_{\mathbf{A}}(\mathbf{C}, \mathbf{C}'(n))\) .

The elements of \(Z_n(\text{Homgr}_A(C, C'))\) are the graded homomorphisms \(f\) of (descending) degree \(n\) of \(C\) to \(C'\) such that \(d' \circ f = (-1)^n f \circ d\) , that is, the morphisms of complexes of \(C\) to \(C'(n)\) , or again of \(C(p)\) to \(C'(p + n)\) for \(p\) any fixed. These are said to be the morphisms of complexes of (descending) degree \(n\) of \(C\) to \(C'\) ; if \(f, g \in Z_n(\text{Homgr}_A(C, C'))\) and \(s \in \text{Homgr}_A(C, C')_{n+1}\) , then the condition \(g - f = Ds\) means that \(s\) is a homotopy connecting the morphisms \(f\) and \(g\) of \(C\) to \(C'(n)\) , so that \(H_n(\text{Homgr}_A(C, C'))\) is the \(k\) -module of homotopy classes of morphisms of degree (descending) \(n\) of \(C\) to \(C'\) .

Let \(\alpha\in\mathrm{H}_{n}(\mathrm{Homgr}_{\mathrm{A}}(\mathrm{C},\mathrm{C}^{\prime}))\) and \(p\in\mathbf{Z}\) . Let us represent \(\alpha\) by \(f\in\mathbf{Z}_{n}(\mathrm{Homgr}_{\mathrm{A}}(\mathrm{C},\mathrm{C}^{\prime}))\) ; then \(f\) is a morphism of complexes of \(\mathbf{C}\) to \(\mathbf{C}^{\prime}(n)\) , hence \(\mathbf{H}_{p}(f)\) is a homomorphism of \(\mathbf{H}_{p}(\mathbf{C})\) to \(\mathbf{H}_{p}(\mathbf{C}^{\prime}(n))=\mathbf{H}_{p+n}(\mathbf{C}^{\prime})\) ; since \(\mathbf{H}_{p}(f)\) depends only on the homotopy class \(\alpha\) of \(f(\mathbf{X},\mathbf{p.~33,~prop.~3})\) , we deduce a canonical homomorphism of \(k\) -modules

\[
\mathrm{H} _ {n} \left(\operatorname{Homgr} _ {\mathrm{A}} \left(\mathrm{C}, \mathrm{C} ^ {\prime}\right)\right)\rightarrow \operatorname{Hom} _ {\mathrm{A}} \left(\mathrm{H} _ {p} (\mathrm{C}), \mathrm{H} _ {p + n} \left(\mathrm{C} ^ {\prime}\right)\right),
\]

whence a graded k-linear map of degree 0, called canonical

\[
\lambda (\mathrm{C}, \mathrm{C} ^ {\prime}): \mathrm{H} (\operatorname{Homgr} _ {\mathbf {A}} (\mathrm{C}, \mathrm{C} ^ {\prime})) \rightarrow \dot {\operatorname{Homgr}} _ {\mathbf {A}} (\mathrm{H} (\mathrm{C}), \mathrm{H} (\mathrm{C} ^ {\prime})).
\]

The homogeneous components of \(\lambda(C, C')\) will often be denoted:

\[
\lambda^ {n} (\mathrm{C}, \mathrm{C} ^ {\prime}): \mathrm{H} ^ {n} \left(\operatorname{Homgr} _ {\mathrm{A}} \left(\mathrm{C}, \mathrm{C} ^ {\prime}\right)\right)\rightarrow \prod_ {p + q = n} \operatorname{Hom} _ {\mathrm{A}} \left(\mathrm{H} _ {p} (\mathrm{C}), \mathrm{H} ^ {q} \left(\mathrm{C} ^ {\prime}\right)\right).
\]

PROPOSITION 1. --- If C is zero on the right and C' zero on the left, then Homgr \(_{A}\) (C, C') is zero on the left, and the canonical k-linear map

\[
\lambda^ {0} (\mathrm{C}, \mathrm{C} ^ {\prime}): \mathrm{H} ^ {0} \left(\operatorname{Homgr} _ {\mathrm{A}} (\mathrm{C}, \mathrm{C} ^ {\prime})\right)\rightarrow \operatorname{Hom} _ {\mathrm{A}} \left(\mathrm{H} _ {0} (\mathrm{C}), \mathrm{H} ^ {0} \left(\mathrm{C} ^ {\prime}\right)\right)
\]

is bijective.

We have exact sequences

\[
0 \longrightarrow \mathrm{H} ^ {0} (\mathrm{C} ^ {\prime}) \stackrel {i} {\longrightarrow} \mathrm{C} ^ {\prime 0} \stackrel {d ^ {\prime 0}} {\longrightarrow} \mathrm{C} ^ {\prime 1}
\]

\[
\mathrm{C} _ {1} \xrightarrow {d _ {1}} \mathrm{C} _ {0} \xrightarrow {p} \mathrm{H} _ {0} (\mathrm{C}) \longrightarrow 0.
\]

On the other hand \(\mathrm{Homgr}_{\mathbf{A}}^{0}(\mathbf{C},\mathbf{C}^{\prime})\) is identified with \(\mathrm{Hom}_{\mathbf{A}}(\mathbf{C}_{0},\mathbf{C}^{\prime 0})\) , \(\mathbf{Z}^{0}(\mathrm{Homgr}_{\mathbf{A}}(\mathbf{C},\mathbf{C}^{\prime}))\) then being identified with the set of \(f:\mathbf{C}_{0}\to \mathbf{C}^{\prime 0}\) such that \(d^{\prime 0}\circ f = 0\) , \(f\circ d_{1} = 0\) ; \(\mathbf{B}^{0}(\mathrm{Homgr}_{\mathbf{A}}(\mathbf{C},\mathbf{C}^{\prime}))\) is zero; finally the map \(\lambda^0\) associates to the class of \(f\) modulo \(\{0\}\) the homomorphism \(\varphi :\mathrm{H}_0(\mathrm{C})\to \mathrm{H}^0 (\mathrm{C}')\) such that \(f = i\circ \varphi \circ p\) , whence the proposition.

Let \(u: \tilde{\mathbf{C}} \to \mathbf{C}\) and \(u': \mathbf{C}' \to \tilde{\mathbf{C}}'\) be morphisms of complexes; then the canonical homomorphism \(\operatorname{Homgr}_{\mathbf{A}}(u, u') : \operatorname{Homgr}_{\mathbf{A}}(\mathbf{C}, \mathbf{C}') \to \operatorname{Homgr}_{\mathbf{A}}(\tilde{\mathbf{C}}, \tilde{\mathbf{C}}')\) , defined by \(f\mapsto u^{\prime}\circ f\circ u\) , is a morphism of complexes, as follows immediately from formula (1). Moreover, the following diagram is commutative

\[
\begin{array}{c} \mathrm{H} (\mathrm{Homgr} _ {\mathsf {A}} (\mathrm{C}, \mathrm{C} ^ {\prime})) \xrightarrow {\lambda (\mathrm{C} , \mathrm{C} ^ {\prime})} \mathrm{Homgr} _ {\mathsf {A}} (\mathrm{H} (\mathrm{C}), \mathrm{H} (\mathrm{C} ^ {\prime})) \\ \mathrm{H} (\mathrm{Homgr} _ {\mathsf {A}} (u, u ^ {\prime})) \Big \downarrow \\ \mathrm{H} (\mathrm{Homgr} _ {\mathsf {A}} (\tilde {\mathsf {C}}, \tilde {\mathsf {C}} ^ {\prime})) \xrightarrow {\lambda (\tilde {\mathsf {C}} , \tilde {\mathsf {C}} ^ {\prime})} \mathrm{Homgr} _ {\mathsf {A}} (\mathrm{H} (\tilde {\mathsf {C}}), \mathrm{H} (\tilde {\mathsf {C}} ^ {\prime}))  . \end{array}
\]

PROPOSITION 2. --- a) Let \(C' \xrightarrow{u} C \xrightarrow{v} C''\) be an exact sequence of A-complexes, \(P\) a projective complex, E an injective complex (X, p. 25). Then the sequences

\[
\operatorname{Homgr} _ {\mathbf {A}} (\mathrm{P}, \mathrm{C} ^ {\prime}) \xrightarrow {\operatorname{Homgr} (1 , u)} \operatorname{Homgr} _ {\mathbf {A}} (\mathrm{P}, \mathrm{C}) \xrightarrow {\operatorname{Homgr} (1 , v)} \operatorname{Homgr} _ {\mathbf {A}} (\mathrm{P}, \mathrm{C} ^ {\prime \prime})
\]

and

\[
\operatorname{Homgr} _ {\mathbf {A}} \left(\mathrm{C} ^ {\prime \prime}, \mathrm{E}\right) \xrightarrow {\operatorname{Homgr} (v , 1)} \operatorname{Homgr} _ {\mathbf {A}} (\mathrm{C}, \mathrm{E}) \xrightarrow {\operatorname{Homgr} (u , 1)} \operatorname{Homgr} _ {\mathbf {A}} \left(\mathrm{C} ^ {\prime}, \mathrm{E}\right)
\]

are exact sequences of complexes of k-modules.

\begin{enumerate}
\def\labelenumi{\alph{enumi})}
\setcounter{enumi}{1}
\tightlist
\item
  Let \(0 \to \mathbf{C}' \xrightarrow{u} \mathbf{C} \xrightarrow{v} \mathbf{C}'' \to 0\) be a sequence of A-complexes which is split as an exact sequence of graded A-modules (this is the case, for example, if \(\mathbf{C}'\) is injective, or if \(\mathbf{C}''\) is projective). Then for every complex E, the sequences
\end{enumerate}

\[
0 \rightarrow \operatorname{Homgr} _ {\mathbf {A}} (\mathrm{E}, \mathrm{C} ^ {\prime}) \xrightarrow {\operatorname{Homgr} (1 , u)} \operatorname{Homgr} _ {\mathbf {A}} (\mathrm{E}, \mathrm{C}) \xrightarrow {\operatorname{Homgr} (1 , v)} \operatorname{Homgr} _ {\mathbf {A}} (\mathrm{E}, \mathrm{C} ^ {\prime \prime}) \rightarrow 0
\]

\[
0 \rightarrow \operatorname{Homgr} _ {\mathbf {A}} \left(\mathrm{C} ^ {\prime \prime}, \mathrm{E}\right) \xrightarrow {\operatorname{Homgr} (v , 1)} \operatorname{Homgr} _ {\mathbf {A}} (\mathrm{C}, \mathrm{E}) \xrightarrow {\operatorname{Homgr} (u , 1)} \operatorname{Homgr} _ {\mathbf {A}} \left(\mathrm{C} ^ {\prime}, \mathrm{E}\right)\rightarrow 0
\]

are exact sequences of complexes of k-modules.

In the case \(a\) ), we note that the sequences

\[
\operatorname{Hom} _ {\mathbf {A}} \left(\mathrm{P} _ {p}, \mathrm{C} _ {q} ^ {\prime}\right)\rightarrow \operatorname{Hom} _ {\mathbf {A}} \left(\mathrm{P} _ {p}, \mathrm{C} _ {q}\right)\rightarrow \operatorname{Hom} _ {\mathbf {A}} \left(\mathrm{P} _ {p}, \mathrm{C} _ {q} ^ {\prime \prime}\right)
\]

and

\[
\operatorname{Hom} _ {\mathbf {A}} \left(\mathrm{C} _ {q} ^ {\prime \prime}, \mathrm{E} _ {p}\right)\rightarrow \operatorname{Hom} _ {\mathbf {A}} \left(\mathrm{C} _ {q}, \mathrm{E} _ {p}\right)\rightarrow \operatorname{Hom} _ {\mathbf {A}} \left(\mathrm{C} _ {q} ^ {\prime}, \mathrm{E} _ {p}\right)
\]

are exact for all \(p, q \in \mathbf{Z}\) , and we apply II, p. 10, prop. 5 and II, p. 13, prop. 7. The proof of \(b\) ) is analogous.

